We study the favourite--longshot bias on event-contract exchanges with three stratified random
samples: quotes on 19,007 settled non-financial Kalshi markets from 2021 to 2026, 2.7 million
Kalshi trade prints, and 2.6 million fills by 237,569 identified Polymarket wallets. On Kalshi,
contracts priced at 10 cents or less return $-55\%$ per dollar at the mid-quote and $-75\%$ after
the spread and fee, while favourites roughly break even. The calibration slope of outcomes on prices
is 1.13 ($t=6.5$ against one). The bias appears in every category and does not fade as resolution
approaches. It is no weaker in liquid markets and has persisted while the venue grew a
thousand-fold. It is steepest in multi-outcome and deadline contracts. Small orders hold the
longshots: 39\% of contracts in taker orders of \$10 or less are longshots, against 1\% in orders
above \$1,000. On Polymarket, retail wallets account for 86\% of net longshot buying. The one-cent
tick, the spread and the fee bound the bias, so taking the other side as a taker does not pay. But
they do not explain it: it survives away from the tick floor, at one-hour horizons and in the most
liquid markets. Insider trading and belief dispersion fit poorly, while one mild
probability-weighting function fits every category. The evidence points to lottery demand by
small traders, met at prices that just cover the cost of supplying it.
